\section{Public monetary histories and the boundary of an empirical claim}
\label{s6:public-dynamic}
The archived statsmodels macroeconomic extract contains 203 quarterly
observations from 1959Q1 through 2009Q3, assembled from FRED and labor
statistics with source access dated 15 December 2009
\citep{public_macrodata_dynamic}. It includes an end-quarter M1 stock,
a quarterly average of monthly three-month Treasury-bill rates, output
and inflation. Its data, source notes and exact numerical implementation
are preserved. This is a real dynamic financial-data analysis with a known
historical snapshot; it is not a current bank deposit-offer dataset.

The fixed projection response is
$100(\log M1_{t+h}-\log M1_{t-1})$ per percentage-point change in
the quarterly-average bill rate. Two lags each of the rate change, money
growth, output growth and inflation enter every regression. Four
specifications are fixed before fitting: these lagged controls; adding
current output growth and inflation; a pre-1984 origin window; and a
post-1984 origin window. Horizons zero through four use common origins
within each specification. Bartlett lag four covariance is computed from
the complete OLS influence, including controls. An independent statsmodels
OLS/HAC call agrees in coefficient and covariance to below $3\times10^{-14}$.

At horizon four, the full-sample lag-control estimate is $-2.102$
percentage points of cumulative money growth, standard error .452.
Adding current output/inflation gives $-2.163$ (.453). The pre-1984
estimate is $-.997$ (.262), while the later window gives $-4.329$ (.685).
There are 196 common full-sample origins and 97/99 in the two windows.
The different windows define different projection populations; selecting
one after observing its significance would add a selection problem. Their
contrast is a registered sensitivity analysis, not proof of a known break
date or a homogeneous policy effect.

These numbers do not identify a bank's price intervention. A quarterly
average rate and end-quarter stock do not supply a within-quarter decision
clock; central-bank information and money demand can confound the
innovation; M1 is not one bank's retained deposits; structural and
measurement regimes can change over fifty years. Current output controls
may be post-policy variables under one causal ordering and predecision
information under another. The data alone cannot settle that ordering.
The inference reported is conditional projection uncertainty, not uncertainty
over all these causal interpretations.

The public-data investigation also inspects the Federal Reserve H.8
description \citep{fed_h8_dynamic}. Its aggregate bank balance-sheet
statistics are useful for surveillance, but do not reveal customer offers,
the private assignment information, first-available entity-linked histories
and representative audited matches required by the treasury brief. Neither
that feed nor the archived macro series identifies the stipulated bank
experiment. Obtaining a current bank-specific price/offer history and its
assignment evidence is therefore a genuine empirical requirement, rather
than an excuse to call any public time series causal.

The central known-truth evidence uses the realistic but explicitly synthetic
funding mechanism, common shocks and revision/linkage construction. At 360
and 1,440 dates, 400 repetitions give financial interval coverage .945 and
.9375; pretending the branches supply independent calendar information
gives .230 and .2175. The corresponding joint Bonferroni response coverage
is .9425 and .965. Binomial Monte Carlo error is approximately one percentage
point for coverage near .95 at 400 repetitions; these differences do not
certify exact nominal calibration. The pair-linkage and audited-servicing
experiments add separate observation and information failures. Taken
together, the public analysis challenges economic interpretation, while
the synthetic analysis checks statistical behavior when its assumptions
are known. Neither estimates a named banking group's causal benefit.

These studies favor different remedies for different failures. Correcting the repricing account changes the economic target; retaining cumulative policy weights changes the dynamic comparison; fitting the private action law changes the observed likelihood; and preserving shared calendars changes uncertainty. Public estimates add sensitivity to interpretation and population, while generated laws expose errors against known truth. The next chapter develops those reliability questions when customer identity, audits, calendars and fitted models are uncertain together.
